% GENERATED FILE. Edit canonical lessons, then run npm run generate:books.
% canonical-source-hash: fnv1a64:45b44c0ad22f876a
% canonical-lessons: MR-C241-kundi, MR-C241-rop, MR-C241-khod, MR-C241-hirval, MR-C241-payvat

\chapter{Flowerpot, Seedling, Trunk, Lawn, Footpath}
\label{ch:mr-flowerpot-seedling-trunk-lawn-footpath}
\hlchaptermodality{241}

\begin{chapteropening}
\textbf{By the end of this chapter:} \emph{I can say a flowerpot, a seedling, a trunk, a lawn and a footpath.}

Complete the last lesson of chapter 241: I can say a flowerpot, a seedling, a trunk, a lawn and a footpath.
\end{chapteropening}

\section[\texorpdfstring{\mr{कुंडी} (kuṇḍī)}{कुंडी (kuṇḍī)}]{\mr{कुंडी} (kuṇḍī) — a flowerpot}

\label{lesson:MR-C241-kundi}



\begin{quote}
Before the new one: say the Marathi for a seat, then the Marathi for a wheel.
\end{quote}

\subsection*{You'll want to know: \mr{कुंडी}}
\textbf{\mr{कुंडी}} — \emph{kuṇḍī} — \textquotedblleft{}a flowerpot\textquotedblright{}.

\begin{grammarlens}[title={What you've built}]
One more word for this chapter.
\end{grammarlens}

\subsection*{Your turn}
\begin{itemize}
\raggedright
  \item \emph{Say it:} \emph{kuṇḍī}
  \item \emph{Say it:} \emph{kuṇḍī}, once more
  \item \emph{Say it:} say \emph{kuṇḍī}
  \item \emph{Recall:} say the Marathi for a seat, then the Marathi for a wheel, then say \emph{kuṇḍī} again
\end{itemize}

\begin{tcolorbox}[breakable,colback=teal!4,colframe=teal!35!black,title={Before you move on}]
What does \mr{कुंडी} mean? (\textquotedblleft{}A flowerpot\textquotedblright{}.) Say it once more.
\end{tcolorbox}
\section[\texorpdfstring{\mr{रोप} (rop)}{रोप (rop)}]{\mr{रोप} (rop) — a seedling, a young plant}

\label{lesson:MR-C241-rop}



\begin{quote}
Before the new one: say the Marathi for a wheel, then the Marathi for a flowerpot.
\end{quote}

\subsection*{You'll want to know: \mr{रोप}}
\textbf{\mr{रोप}} — \emph{rop} — \textquotedblleft{}a seedling, a young plant\textquotedblright{}.

\begin{grammarlens}[title={What you've built}]
One more word for this chapter.
\end{grammarlens}

\subsection*{Your turn}
\begin{itemize}
\raggedright
  \item \emph{Say it:} \emph{rop}
  \item \emph{Say it:} \emph{rop}, once more
  \item \emph{Say it:} say \emph{rop}
  \item \emph{Recall:} say the Marathi for a wheel, then the Marathi for a flowerpot, then say \emph{rop} again
\end{itemize}

\begin{tcolorbox}[breakable,colback=teal!4,colframe=teal!35!black,title={Before you move on}]
What does \mr{रोप} mean? (\textquotedblleft{}A seedling, a young plant\textquotedblright{}.) Say it once more.
\end{tcolorbox}
\section[\texorpdfstring{\mr{खोड} (khoḍ)}{खोड (khoḍ)}]{\mr{खोड} (khoḍ) — a trunk}

\label{lesson:MR-C241-khod}



\begin{quote}
Before the new one: say the Marathi for a flowerpot, then the Marathi for a seedling.
\end{quote}

\subsection*{You'll want to know: \mr{खोड}}
\textbf{\mr{खोड}} — \emph{khoḍ} — \textquotedblleft{}a trunk\textquotedblright{}.

\begin{grammarlens}[title={What you've built}]
One more word for this chapter.
\end{grammarlens}

\subsection*{Your turn}
\begin{itemize}
\raggedright
  \item \emph{Say it:} \emph{khoḍ}
  \item \emph{Say it:} \emph{khoḍ}, once more
  \item \emph{Say it:} say \emph{khoḍ}
  \item \emph{Recall:} say the Marathi for a flowerpot, then the Marathi for a seedling, then say \emph{khoḍ} again
\end{itemize}

\begin{tcolorbox}[breakable,colback=teal!4,colframe=teal!35!black,title={Before you move on}]
What does \mr{खोड} mean? (\textquotedblleft{}A trunk\textquotedblright{}.) Say it once more.
\end{tcolorbox}
\section[\texorpdfstring{\mr{हिरवळ} (hirvaḷ)}{हिरवळ (hirvaḷ)}]{\mr{हिरवळ} (hirvaḷ) — a lawn}

\label{lesson:MR-C241-hirval}



\begin{quote}
Before the new one: say the Marathi for a seedling, then the Marathi for a trunk.
\end{quote}

\subsection*{You'll want to know: \mr{हिरवळ}}
\textbf{\mr{हिरवळ}} — \emph{hirvaḷ} — \textquotedblleft{}a lawn\textquotedblright{}.

\begin{grammarlens}[title={What you've built}]
One more word for this chapter.
\end{grammarlens}

\subsection*{Your turn}
\begin{itemize}
\raggedright
  \item \emph{Say it:} \emph{hirvaḷ}
  \item \emph{Say it:} \emph{hirvaḷ}, once more
  \item \emph{Say it:} say \emph{hirvaḷ}
  \item \emph{Recall:} say the Marathi for a seedling, then the Marathi for a trunk, then say \emph{hirvaḷ} again
\end{itemize}

\begin{tcolorbox}[breakable,colback=teal!4,colframe=teal!35!black,title={Before you move on}]
What does \mr{हिरवळ} mean? (\textquotedblleft{}A lawn\textquotedblright{}.) Say it once more.
\end{tcolorbox}
\section[\texorpdfstring{\mr{पायवाट} (pāyvāṭ)}{पायवाट (pāyvāṭ)}]{\mr{पायवाट} (pāyvāṭ) — a footpath}

\label{lesson:MR-C241-payvat}



\begin{quote}
Before the new one: say the Marathi for a trunk, then the Marathi for a lawn.
\end{quote}

\subsection*{You'll want to know: \mr{पायवाट}}
\textbf{\mr{पायवाट}} — \emph{pāyvāṭ} — \textquotedblleft{}a footpath\textquotedblright{}.

\begin{grammarlens}[title={What you've built}]
One more word for this chapter.
\end{grammarlens}

\subsection*{Your turn}
\begin{itemize}
\raggedright
  \item \emph{Say it:} \emph{pāyvāṭ}
  \item \emph{Say it:} \emph{pāyvāṭ}, once more
  \item \emph{Say it:} say \emph{pāyvāṭ}
  \item \emph{Recall:} say the Marathi for a trunk, then the Marathi for a lawn, then say \emph{pāyvāṭ} again
\end{itemize}

\begin{tcolorbox}[breakable,colback=teal!4,colframe=teal!35!black,title={Before you move on}]
What does \mr{पायवाट} mean? (\textquotedblleft{}A footpath\textquotedblright{}.) Say it once more.
\end{tcolorbox}
